\chapter{Discussion and Conclusion}\label{sec:discussion}

\section{Summary of contributions}
This work presents an integrated workflow for constructing, representing, training on, and validating route-aware ETA-oriented traffic datasets. The contribution is not a single model, a single screenshot-driven application, or a single exported dataset. It is the connection between these parts: a desktop construction environment, a shared graph representation, learning-based validation, and a web loop that turns model predictions into persistent evidence.

The first contribution is the desktop construction workflow. It supports two different sources of traffic data: controlled SUMO simulation and repaired GPS trajectory conversion. The simulation path provides a configurable environment with zones, vehicle types, schedules, landmarks, and sampled traffic state. The trajectory path supports a less controlled but more realistic source by preparing an OSM-derived network, repairing noisy trajectories, matching them to SUMO routes, and exporting them into the same downstream format.

The second contribution is the shared route-aware dynamic graph representation described in Chapter~\ref{sec:dataset}. Instead of exporting disconnected trip tables, the workflow preserves road topology, vehicle state, remaining route context, dynamic edge relations, and ETA labels in a time-indexed graph structure. This representation is deliberately shaped for ETA learning, where the route a vehicle intends to follow is central to the prediction task.

The third contribution is validation across two levels. Chapter~\ref{sec:eval} shows that the representation can train a DSTRA-compatible ETA model on both construction paths. Chapter~\ref{sec:web} shows that the trained model can then be exercised in the \WebApp{}, where selected journeys produce logged prediction records with realized outcomes and aggregate analysis. Together, these chapters show that the dataset is not only generated successfully, but also usable by a neural ETA model and by an interactive evaluation workflow.

\section{What the shared representation achieved}
The key design question was whether simulation-derived traffic and trajectory-derived traffic could be expressed through one model-facing contract. These sources differ substantially. Simulation provides strong observability and controlled demand; trajectory conversion begins with sparse, noisy real-world traces and must infer route structure before graph snapshots can be emitted. The value of the representation is that it absorbs this difference at construction time, so downstream consumers do not need separate schemas or special-purpose adapters.

The invariants introduced in Chapter~\ref{sec:dataset} are what make this possible. The static substrate keeps road and junction identifiers stable across all snapshots. The dynamic layer records time-indexed vehicle state and road-level traffic descriptors. Route conditioning attaches planned or inferred route context to each vehicle, and labels remain separated from features so that supervised learning targets do not leak into the input graph. These choices make the exported bundles comparable even when their construction paths are different.

This path independence is important for reproducibility. A researcher can train or test a model against a simulation bundle and a trajectory-conversion bundle without rewriting preprocessing logic. A developer can also inspect the same conceptual objects in both cases: roads, junctions, vehicles, dynamic relations, snapshots, and labels. The result is a workflow in which construction details remain visible and auditable, but the learning and web-evaluation layers see a common interface.

\section{Interpretation of learning validation}
The learning results in Chapter~\ref{sec:eval} support the claim that the representation is useful as a model input, not only as a storage format. DSTRA converges on both the simulation and trajectory-conversion bundles, and in both cases it substantially outperforms a naive mean-duration baseline. This indicates that the graph contains predictive signal beyond the marginal distribution of trip durations.

The simulation path converges faster and reaches slightly lower error. This is expected because its snapshots are denser in active traffic and contain more frequent dynamic vehicle and relation edges. Those repeated interaction patterns make it easier for the model to learn how evolving traffic state affects ETA. The trajectory-conversion path remains harder: sparse GPS observations must be repaired, split when invalid segments occur, matched to road-network routes, and projected into graph snapshots. Even with that additional uncertainty, the model still converges, which is the central validation point for the conversion branch.

It is also important to interpret the model role correctly. The report does not claim DSTRA as the novelty of this project. DSTRA is used as a demanding downstream learner because it depends on exactly the information the representation aims to preserve: dynamic graph snapshots, route context, vehicle state, and ETA labels. Its successful training is therefore evidence that the dataset format supports neural ETA modeling.

\section{Interpretation of web validation}
Chapter~\ref{sec:web} extends validation from offline learning to application-level use. The web loop checks whether a trained model can be used in an end-to-end workflow: load a compatible bundle, select a route, invoke ETA inference, execute or replay the journey, record the realized outcome, and aggregate the error evidence. This matters because a model checkpoint alone does not prove that the surrounding toolchain is usable.

The close agreement between training-evaluation metrics and web-tool journey measurements is especially important. If the web-tool measurements had diverged strongly from the offline results, that would suggest a mismatch between the model evaluation protocol and the deployed interaction loop. Instead, the measured errors are similar, which supports the validity of the logged web results and reinforces that the generated datasets can be used to train neural-network ETA models that remain usable in the application.

The persistence layer also changes the nature of the evidence. Prediction results are not temporary screen outputs; they become stored journey records that can be revisited, filtered, plotted, and reported. This gives the workflow a practical audit trail from dataset construction through model use, and it makes later analysis possible without rerunning every journey.

\section{Limitations and threats to validity}
The main limitation is dataset scope. The report uses representative bundled examples rather than an exhaustive set of cities, networks, demand patterns, and trajectory sources. The results therefore demonstrate feasibility and workflow coherence, but they should not be read as universal performance claims for all traffic environments.

Simulation realism is another limitation. SUMO provides control, repeatability, and complete observability, which are valuable for constructing benchmark-like data. However, controlled simulation may omit behavioral complexity, sensing noise, routing heterogeneity, and unexpected disruptions found in real traffic. Good performance on simulation-derived data is therefore evidence of representation usability, not proof of real-world deployment readiness.

The trajectory-conversion path has different risks. GPS traces can be sparse, repeated points can hide true trip starts or ends, and unrealistic jumps can create invalid segments. Route matching may choose plausible but incorrect road paths, especially in dense urban networks or when the source trajectory is noisy. These uncertainties can propagate into graph snapshots and labels, making the converted dataset harder to learn from and harder to interpret.

The model scope is also intentionally limited. DSTRA is one compatible learner, not a comprehensive benchmark over all ETA models. Other graph neural networks, sequence models, tree-based methods, or hybrid route models may respond differently to the same representation. The present results show that the representation can support a neural model; they do not establish that the chosen model is optimal.

Three additional threats deserve explicit mention. First, \textbf{label leakage}: the representation keeps ETA labels in separate label files, physically decoupled from the snapshot feature tensors. This structural separation prevents the target from contaminating the graph input during training, but readers implementing their own loaders should verify that their data pipeline does not inadvertently read from label files at feature-extraction time. Second, \textbf{sensitivity to map-matching errors}: when the route-matching algorithm selects an incorrect road path, the error propagates silently into route-left features and dynamic edge relations for the entire trajectory segment. The trajectory-conversion results therefore represent an upper bound on achievable accuracy under noisy matching, and the true effect of matching errors on learning is not separately quantified here. Third, \textbf{limited external validity}: all experiments use one simulated three-zone network and one Portuguese urban taxi dataset. It is unknown how well the representation, the matching procedure, or the trained model would transfer to other cities, road network styles, transit modes, or demand patterns. Claims about generalization across geographies should await experiments on additional benchmark networks.

The choice of evaluation metrics also carries a limitation. MAPE is sensitive to short trip durations: a small absolute error on a two-minute journey produces a high percentage error that inflates the aggregate. Tables~\ref{tab:learning_validation_results} and~\ref{tab:web_eval_results} therefore report MedAE and P90/P95 absolute errors alongside MAE, RMSE, and MAPE. The MedAE values confirm that the median prediction error is well below the mean, and the P90/P95 columns bound the tail, giving a more complete picture of model behavior on short-trip populations.

Finally, the web tool is a validation harness rather than a production navigation system. It demonstrates route selection, ETA inference, journey logging, aggregate plots, and report-style evidence. It does not address production requirements such as live map updates, user-scale concurrency, robustness to service outages, calibrated uncertainty, privacy controls, or deployment across arbitrary cities without preparation. Aggregate metrics such as MAE, RMSE, and MAPE also do not capture every operational concern, including regional fairness, rare-event behavior, or prediction stability under disruptions.

\section{Reproducibility}
The dataset generator, trajectory-conversion tool, web evaluation client, and this report are all available in the public repository \href{https://github.com/Turgibot/Multi-Variant-Simulated-Traffic-Dataset-Creator-and-Model-Tester}{Turgibot/Multi-Variant-Simulated-Traffic-Dataset-Creator-and-Model-Tester}. The pre-generated graph datasets used in this report are distributed separately in a \href{https://drive.google.com/drive/folders/1baEONbv8WyDVanF_qvLCq1UJ4oT5N7LE?usp=sharing}{shared dataset folder}. The web evaluation application is deployed at \href{https://smarttransportationlab.com/sim-demo}{smarttransportationlab.com/sim-demo} and its source is available at \href{https://github.com/Turgibot/TrafficLab}{Turgibot/TrafficLab}.

\paragraph{Software versions.}
\begin{sloppypar}
The toolchain was developed and tested with SUMO~1.20, Python~3.10, PyTorch~2.2, and PyTorch Geometric~2.5. The web backend runs FastAPI~0.111 with Python~3.10; the frontend uses Vue.js~3.
\end{sloppypar}

\paragraph{Hardware.}
Training experiments were run on a single NVIDIA GPU (16\,GB VRAM). Dataset generation and trajectory conversion run on CPU and do not require a GPU.

\paragraph{Random seeds.}
The DSTRA training runs used a fixed random seed applied to both PyTorch and NumPy to ensure reproducibility of weight initialization and data shuffle.
The seed value is recorded in the training configuration file committed to the repository.

\paragraph{Reproducing evaluation tables and figures.}
\begin{enumerate}
    \item Clone the repository and follow the \texttt{README} to install dependencies.
    \item Download the pre-generated datasets from the shared dataset folder.
    \item Run the training script with the bundled configuration to reproduce the held-out metrics in Table~\ref{tab:learning_validation_results}.
    \item Load the exported bundles into the web client and run evaluation sessions to reproduce the web-tool measurements in Table~\ref{tab:web_eval_results}.
    \item Scatter plots and training curves are generated automatically by the evaluation and training scripts; outputs are written to the \texttt{results/} directory.
\end{enumerate}
To reproduce this report PDF, run \texttt{./compile.sh} from \texttt{Academia/Project\_Report/}; the script resolves the full \texttt{pdflatex} and \texttt{bibtex} chain automatically.

\section{Data licensing and privacy}
This work uses two external data sources, each with distinct licensing conditions.

\paragraph{Porto taxi dataset.}
The Porto taxi trajectory data was released as part of the ECML/PKDD 2015 taxi-service trajectory prediction competition~\cite{porto2015pkdd} and is publicly available on Kaggle for research use. The dataset contains GPS traces collected from taxi vehicles operating in Porto, Portugal; it does not include direct personally identifiable information such as names or payment records. However, trajectory data can be re-identifying in practice: sufficiently long or distinctive routes can be linked to individuals even without explicit identifiers. The system uses this dataset strictly as a conversion demonstrator. Any use of the Porto data beyond the competition's original research license, or in contexts where re-identification risk is relevant, requires compliance with applicable privacy regulations.

\paragraph{OpenStreetMap data.}
Road network geometry for trajectory conversion and simulation scenarios is obtained from OpenStreetMap (OSM) via the Overpass API~\cite{osm}. OSM data is distributed under the Open Database License (ODbL), which requires attribution and share-alike conditions on derived databases. Users of this toolchain who export or redistribute network files derived from OSM data must comply with the ODbL, including providing appropriate credit to OpenStreetMap contributors.

\paragraph{General principle.}
The toolchain itself does not collect, store, or transmit personal data. The web evaluation client stores journey records (route, predicted and realized travel times, and aggregate statistics) in a local database for analysis purposes; these records do not contain vehicle identifiers or user identities. Any deployment of the web evaluation client in a context involving real users or real vehicles should be assessed against applicable data-protection regulations before production use.

\section{Future work}
Future work should broaden the empirical base of the framework. Additional simulation networks, demand schedules, city scales, and trajectory datasets would make it possible to separate representation effects from scenario-specific behavior. More systematic ablations could also test which features are most important: route suffixes, dynamic relation edges, road-demand features, temporal history, or compression choices.

The trajectory-conversion branch would benefit from stronger validation and uncertainty tracking. Route-matching confidence scores, explicit labels for repaired segments, and diagnostics for ambiguous paths could help downstream users understand which converted examples are most reliable. This would also support training regimes that weight examples by confidence or exclude low-quality matches from sensitive evaluations.

The modeling layer can be expanded beyond DSTRA. Since the representation is exported as a reusable graph artifact, future work can compare additional neural and non-neural predictors under the same construction assumptions. This would turn the toolchain into a stronger benchmark environment rather than a validation workflow for one bundled model.

The web environment can also grow into a richer analysis platform. Useful extensions include more cohort analyses by time of day, trip length, origin/destination zones, vehicle type, and error drift across sessions. Report generation could become more automated, and stored journeys could be linked back to the exact dataset configuration, model checkpoint, and inference settings used during prediction.

Finally, packaging more reference scenarios would make the work easier to compare externally. A set of prepared networks with documented demand patterns, noise levels, and trajectory-conversion settings would let other researchers evaluate models under shared construction assumptions instead of relying on opaque CSV extracts or one-off preprocessing scripts.

\section{Conclusion}
Overall, the project shows a reusable route-aware ETA dataset workflow in which controlled simulation and real trajectory conversion both produce learnable, web-testable graph data. The desktop tool makes construction explicit, the shared representation provides a stable model-facing contract, DSTRA training validates learnability, and the web client turns model predictions into persistent evaluation evidence. This integrated chain is the central contribution: it narrows the gap between dataset generation, neural ETA modeling, and practical validation.
